\documentclass{article}
\usepackage{enumitem}
\usepackage{amsmath}
\begin{document}

\title{Chapter 11: Transport in Plants}
\date{}
\maketitle

\begin{enumerate}[label=Q\arabic*.]

\item The driving force for water movement up a tall tree in the xylem is primarily:
\\(A) Root pressure
\\(B) Transpiration pull (cohesion-tension theory)
\\(C) Osmotic pressure of phloem sap
\\(D) Capillary action alone

\item Water potential ($\Psi_w$) is the sum of:
\\(A) Osmotic potential and turgor pressure only
\\(B) Solute potential ($\Psi_s$) and pressure potential ($\Psi_p$)
\\(C) Matric potential and osmotic potential only
\\(D) Gravitational potential and osmotic potential

\item In plasmolysis, which of the following correctly describes what happens?
\\(A) Cell gains water; protoplast pushes against cell wall
\\(B) Cell loses water; protoplast shrinks away from cell wall
\\(C) Cell bursts due to excess turgor
\\(D) Cell wall dissolves in concentrated solution

\item The Casparian strip in the endodermis forces water to take which pathway into the vascular cylinder?
\\(A) Apoplastic pathway
\\(B) Symplastic pathway (through cytoplasm)
\\(C) Transmembrane pathway
\\(D) Both apoplastic and symplastic equally

\item Guttation is caused by:
\\(A) Transpiration from leaves during day
\\(B) Root pressure forcing water out through hydathodes at night
\\(C) Evaporation from stomata during cold nights
\\(D) Phloem exudation through leaf margins

\item Assertion: Active transport requires ATP and can move solutes against concentration gradients.\\
Reason: Carrier proteins in membranes use energy from ATP hydrolysis to pump solutes.
\\(A) Both true, Reason is correct explanation
\\(B) Both true, Reason is not correct explanation
\\(C) Assertion true, Reason false
\\(D) Assertion false, Reason true

\item Translocation in phloem is explained by the pressure flow hypothesis. According to this model, flow is from:
\\(A) Source (high solute, high pressure) to sink (low solute, low pressure)
\\(B) Sink to source
\\(C) Root to leaf only
\\(D) Leaf to root only

\item Which mineral is part of the chlorophyll molecule?
\\(A) Iron
\\(B) Manganese
\\(C) Magnesium
\\(D) Calcium

\item Stomatal opening is triggered by:
\\(A) High CO$_2$ concentration inside guard cells
\\(B) Active transport of K$^+$ ions into guard cells, increasing osmotic pressure
\\(C) Loss of turgor in guard cells
\\(D) Accumulation of ABA in guard cells

\item Which of the following best describes the apoplastic pathway of water movement?
\\(A) Water moves through the cytoplasm and plasmodesmata
\\(B) Water moves through cell walls and intercellular spaces without crossing membranes
\\(C) Water moves through tonoplast into vacuoles
\\(D) Water moves via aquaporins through membranes

\item Aquaporins facilitate:
\\(A) Active transport of ions
\\(B) Rapid osmotic transport of water across membranes
\\(C) Transport of sugars through phloem
\\(D) Ion exchange at the plasma membrane

\item In the cohesion-tension theory, cohesion refers to:
\\(A) Attraction of water molecules to xylem cell walls
\\(B) Attraction between water molecules due to hydrogen bonding
\\(C) Tension created by transpiration at leaf surfaces
\\(D) Root pressure pushing water upward

\item Wilting of plants occurs when:
\\(A) Turgor pressure exceeds wall pressure
\\(B) Water loss exceeds water absorption, reducing turgor pressure
\\(C) Stomata close fully during night
\\(D) Phloem transport stops

\item Match the following transport mechanisms:\\
1. Simple diffusion $\rightarrow$ (a) No energy, no carrier, high to low concentration\\
2. Facilitated diffusion $\rightarrow$ (b) No energy, needs carrier protein\\
3. Active transport $\rightarrow$ (c) Requires ATP, against gradient\\
4. Osmosis $\rightarrow$ (d) Water movement through semipermeable membrane
\\(A) 1-a, 2-b, 3-c, 4-d
\\(B) 1-b, 2-a, 3-d, 4-c
\\(C) 1-c, 2-d, 3-a, 4-b
\\(D) 1-d, 2-c, 3-b, 4-a

\item Which of the following conditions would increase the rate of transpiration?
\\(A) High humidity, low wind, low temperature
\\(B) Low humidity, high wind, high temperature, and open stomata
\\(C) High CO$_2$ concentration around leaves
\\(D) Water deficit in soil

\end{enumerate}
\end{document}
